\documentclass[11pt]{article}
\usepackage{resolution-paper}

\title{\textbf{Resolution Semantics for Partial Algebras}\\
Intrinsic Finite-Complement Separation and Observational Completion,\\
with a Lean-Checked Many-Sorted Encoding Bridge}
\author{Adrian Puha}
\date{18 August 2026}

\begin{document}
\maketitle

\begin{abstract}
We study one-sorted binary partial algebras through compatible total extensions
that preserve the partial base pointwise while allowing only finitely many new
states outside it.  The underlying normalized Answer algebra is a concrete free
compatible completion; that substrate is classical in character and is not the
novelty claim of the paper.  The main results concern the relative observation
system built over it.

The first structural principle is \emph{relative finite-pattern realization}:
for every finite family of generated Answers there is one compatible
finite-complement observer that is simultaneously injective on that family.
Pairwise finite-tag separation is the two-point consequence, with the explicit
constructor-size bound from the original separator.  These observers induce a
separated filtration and hence a Cauchy completion.

The second structural principle isolates the source of new completion points.
If, at each observation stage, a generated unary family is interpreted as a
deterministic trajectory admitting a uniformly finite trajectory code, then
factorial sampling is Cauchy.  This finite-dynamics step is standard and is not
claimed as new.  The Resolution-specific anti-limit mechanism is complementary:
for a syntactically growing unary orbit, a candidate-tailored finite-pattern
observer realizes the candidate and a long orbit prefix exactly, sends the next
unseen iterate to an absorbing overflow state, and therefore separates every
sufficiently late iterate from the candidate.  Thus trajectory compression plus
finite-pattern escape forces the observational completion to be proper.

Both the exact finite-base properness criterion and the arbitrary-base
base-fixing-context theorem factor through this combined mechanism.  In the
old-fixing case the stage-$n$ indistinguishable factorial pair sharpens to
$c_{(n+1)!}\equiv_n c_{(n+2)!}$, giving unbounded separation ranks.  Natural and
integer arithmetic instantiate the arbitrary-base criterion while retaining
zero-denominator applications as structured Answers rather than ordinary
numbers.  The completed algebra preserves exactly the universal equations in
the language with named base constants.

Finally, a separate Lean-checked many-sorted development formalizes the
operations-only fragment of Diaconescu's $\Gamma$--$\alpha$--$\beta$--$\gamma$
encoding for arbitrary finite arities, including the satisfaction condition,
adjunction laws, semantic consequence, and the fixed-signature initiality
results used here.  Its singleton-sort, no-$TF$, binary-$PF$ specialization is
canonically equivalent to the binary Resolution implementation.  Lean~4.33.0
checks the complete formal package.
\end{abstract}

\tableofcontents

\section{Introduction}\label{sec:introduction}

\subsection{A first example}

Take the natural numbers with $+$, $\cdot$ and $/$, where division is undefined
exactly when the denominator is zero.  We want to total\-ize this partial
algebra without deciding what $0/0$ is.

The standard move is to keep the failed application as syntax.  Applications
that the base algebra can evaluate reduce to base values; the others remain
suspended terms.  So $6/2$ becomes the embedded natural $3$, while
\[
  0/0 \;=\; \susp(\mathsf{div},\old(0),\old(0))
\]
stays a term.  Nothing here is new: this is the free compatible completion, and
we use it only as a substrate.

Now ask a different question.  How much new structure must an observer add in
order to \emph{tell such terms apart}?

Fix a total algebra $T$ together with an injection of the base $\mathbb N$ into
it that fixes every natural pointwise, and require that $T$ add only
\emph{finitely many} states outside that copy of $\mathbb N$.  The base is
infinite, so $T$ is infinite; the budget counts only genuinely new semantic
states.  Call such a $T$ an \emph{observer}.

Two things happen, and they pull in opposite directions.

\begin{itemize}
  \item Given any finite list of terms, one observer separates all of them at
        once.  It keeps $\mathbb N$ pointwise, gives a distinct tag to each
        selected suspended subterm, and sends everything else to a single
        overflow state $\overflow$.
  \item Yet some sequences cannot be followed to a limit.  Start from $0/0$ and
        iterate the context $x\mapsto x+0$:
        \[
          0/0,\qquad 0/0+0,\qquad (0/0+0)+0,\qquad\ldots
        \]
        Because $a+0=a$ holds for every natural, this context fixes the entire
        base pointwise, so no observer can use base elements as memory for it.
        The terms grow syntactically forever, while each observer has only
        finitely many tags.  Given any candidate limit $x$, an observer can
        retain $x$ and a long prefix of the orbit exactly --- and then the next
        iterate must fall into $\overflow$, which is absorbing.  So $x$ is
        separated from every sufficiently late term of the sequence.
\end{itemize}

The sequence is therefore Cauchy for the separation these observers induce, and
it converges to nothing in the algebra.  The completion is \emph{proper}: it
has points that are not terms.  That is the phenomenon this paper is about.

\subsection{The general setting}

Let $D$ be a one-sorted binary partial algebra with carrier $A$.  We consider
compatible total extensions $T$ whose embedding of $A$ is injective, which fix
every element of $A$, and for which $|T\setminus A|$ is finite.  We write
$\overflow$ for the overflow state and $\stageeq{n}$ for indistinguishability
under all observers of budget $n$.

\begin{remark}[On the word \emph{observational}]
We use \emph{observational} in the metric sense: observers induce a separated
filtration and hence a completion.  This is not the observational
(behavioural) satisfaction of algebraic specification, where an observational
equivalence between algebras is induced by a distinguished set of observable
sorts \cite{BidoitHennicker2006,SannellaTarlecki2012}.  No quotient by an
equivalence is taken here; every base element is retained individually.
\end{remark}

The two mechanisms of the example are the two theorems of the paper.

\paragraph{Finite-pattern realization (\cref{thm:finite-pattern-realization}).}
For every finite family of generated Answers there is a single observer
injective on the whole family.  Pairwise separation is the two-point case, not
the primitive one.

\paragraph{Compression and escape (\cref{thm:compression-escape-properness}).}
Two independent facts combine:
\[
  \boxed{
  \text{finite trajectory compression}
  \;+\;
  \text{finite-pattern escape}
  \;\Longrightarrow\;
  \text{proper observational completion}.}
\]
Compression says that if the orbit relevant to a chosen unary family has a
uniformly finite code at each stage, then its factorial sample is Cauchy; this
is pigeonhole periodicity and factorial divisibility.  Escape says that no
generated Answer is a limit of a syntactically growing orbit, by the
candidate-tailored construction of the example.  Compression makes the sample
indistinguishable at every fixed stage; escape stops it converging.  Neither
alone gives properness.

\subsection{Why the observer class is the point}

The reader may expect the finite-family theorem to follow from the pairwise
case by taking products of separators, as in ordinary residual arguments.  It
does not, and the obstruction is visible already at the level of carriers.

Each one-point extension $\mathbb N\hookrightarrow\mathbb N\sqcup\{\overflow\}$
has a one-element complement.  But the diagonal copy of $\mathbb N$ inside the
product of two such extensions has infinitely many points $(a,\overflow)$
outside it.  The class ``fixes the base pointwise, finite complement'' is thus
\emph{not closed under binary products} over an infinite base
(\cref{prop:relative-product-obstruction}).  The finite-family theorem is
therefore proved by direct construction inside the class.

The same restriction is what makes escape non-trivial.  When the base is not
finitely codable, the candidate-tailored observer cannot itself be encoded into
any finite carrier, even though its external complement is finite.  It still
separates the candidate from the whole late tail.

\subsection{Quantitative consequence}

At stage $n$ the old-fixing orbit sees only $n+1$ external states --- the $n$
tags plus $\overflow$ --- unless it enters the pointwise-fixed base, after
which it stabilizes.  Hence for a distinct pair of iterates
\[
  c_{(n+1)!}\stageeq{n}c_{(n+2)!},
\]
and the least separation ranks are unbounded.  The factorial indices only
synchronize finite periods; the content is that finite compression and
finite-pattern escape coexist in the same observer class.

\subsection{Notation}

\begin{center}
\small
\begin{tabular}{@{}ll@{}}
\toprule
$D$, $A$ & partial algebra and its carrier; $A$ may be infinite\\
$\old(a)$ & the embedded copy of a base element $a\in A$\\
$\susp(f,s,t)$ & a suspended application: $f$ was undefined on $s,t$\\
$\RAlg{D}$ & the generated Answer algebra (free compatible completion)\\
observer & compatible total extension fixing $A$ pointwise, finite complement\\
$\overflow$ & the single overflow state of an observer\\
$\stageeq{n}$ & indistinguishable by every observer of budget $n$\\
$\Comp{\RAlg{D}}$ & the completion of the induced filtration\\
proper & the completion has points outside the image of $\RAlg{D}$\\
\bottomrule
\end{tabular}
\end{center}

\subsection{Organization and scope}

\Cref{sec:free-completion} records the substrate.
\Cref{sec:finite-separation} develops the observer class and proves
finite-pattern realization.
\Cref{sec:completion} builds the filtration and proves compression, escape and
their combination; the finite-base and old-fixing criteria follow as instances.
\Cref{sec:equations,sec:arithmetic} give exact equational conservativity and
the arithmetic examples.
\Cref{sec:diaconescu-bridge} locates the construction inside the standard
many-sorted partial-to-total encoding.
\Cref{sec:related-work,sec:limitations} position the work and list what is
open.

\paragraph{Scope.}
Free compatible completion, finite-state periodicity, factorial
synchronization, generic Cauchy completion, and abstract finite-set
discrimination are background.  What is offered here is the relative observer
class over a pointwise-fixed and possibly infinite base, the finite-pattern
realization theorem inside a class that is not product-closed, and the way
escape interacts with trajectory compression to force properness over infinite
bases.  Proofs are given in full in ordinary mathematical language; the
accompanying machine-checked development is described in
\cref{app:formalization} and is not relied on anywhere in the argument.

\section{Generated Answers as a Free Compatible Completion}
\label{sec:free-completion}

This section records the algebraic substrate used by the observational results.
Its role is foundational rather than novel: the normalized construction is a
concrete binary presentation of a free compatible completion.

Let $D$ have carrier $A$, binary operation symbols $\Sigma$, and evaluator
\[
  \eval_D:\Sigma\to A\to A\to\operatorname{Option}(A).
\]
Expressions and raw Answers are generated by
\[
 e::=\operatorname{val}(a)\mid\operatorname{app}(f,e_1,e_2),
 \qquad
 r::=\old(a)\mid\susp(f,r_1,r_2).
\]
Here $A$ is consistently called the \emph{base} in the mathematical prose.
The symbol $\old(a)$ is retained only as the constructor name for the embedded
base leaf in the formal Answer syntax, matching the Lean implementation; it
does not introduce a second mathematical notion distinct from the base.
Define the one-step resolution operation by
\[
\operatorname{liftOp}_D(f,x,y)=
\begin{cases}
 \old(c),&x=\old(a),\ y=\old(b),\ \eval_D(f,a,b)=\operatorname{some}(c),\\
 \susp(f,x,y),&\text{otherwise}.
\end{cases}
\]
The evaluator $\Res_D$ recursively resolves both children and then applies
$\operatorname{liftOp}_D$.  The generated Answer algebra $\RAlg{D}$ is the
image of $\Res_D$, with base embedding $\eta_0(a)=\old(a)$ and operations
induced by $\operatorname{liftOp}_D$.

\begin{proposition}[Conservativity on the base]
\label{prop:generated-conservativity}
If $\eval_D(f,a,b)=\operatorname{some}(c)$, then
\[
 f_{\RAlg{D}}(\eta_0(a),\eta_0(b))=\eta_0(c).
\]
The base embedding $\eta_0$ is injective.
\end{proposition}

A compatible total algebra over $D$ consists of a total $\Sigma$-algebra $T$
and a map $i:A\to T$ satisfying
\[
 \eval_D(f,a,b)=\operatorname{some}(c)
 \Longrightarrow f_T(i(a),i(b))=i(c).
\]
The map need not be injective for the free universal property.

Every raw Answer has a canonical fold into $T$:
\[
 \Phi_T(\old(a))=i(a),\qquad
 \Phi_T(\susp(f,x,y))=f_T(\Phi_T(x),\Phi_T(y)).
\]
Compatibility is exactly what is needed to show that this fold respects the
resolution step.

\begin{theorem}[Free compatible completion]\label{thm:free-completion}
For every compatible total algebra $T$ over $D$, there is a unique homomorphism
\[
 \Phi_T:\RAlg{D}\to T
\]
with $\Phi_T\circ\eta_0=i$.
\end{theorem}

The expression algebra modulo equality after resolution is likewise
canonically equivalent to the generated Answer algebra.

\begin{remark}[Positioning]
The theorem above should be read as a concrete normalized presentation of an
established free-completion phenomenon, not as the main novelty claim.  Its
purpose in this paper is to provide canonical generated syntax and canonical
interpretation maps into compatible observers.  All later separation and
completion statements are relative to those interpretation maps.
\end{remark}

Generated Answers are normal forms: a suspended node cannot have two old
children on which the base evaluator is defined, because such a node would have
reduced.  This simple normalization invariant is the key local fact used by the
finite-pattern observers in the next section.

\section{Relative Finite-Pattern Realization}
\label{sec:finite-separation}

The free-completion theorem gives canonical maps into compatible total
extensions, but it does not say which distinctions between suspended Answers
are visible in independently specified models.  We now restrict attention to
observers that preserve the base pointwise and have only finitely many states
outside it.

Let $T$ be a compatible total extension with injective base embedding
$i:A\hookrightarrow T$.  Its genuine outside part is
\[
 \operatorname{Out}_D(T)=\{z\in T\mid \forall a\in A,\ z\ne i(a)\}.
\]
An intrinsic observer has budget $n$ when
$\operatorname{Out}_D(T)$ injects into $\Fin(n+1)$.  The base $A$ may be
infinite.  The budget counts only states beyond the fixed copy of $A$.

A canonical stage-$n$ representative has carrier
\[
 A\sqcup\Fin(n)\sqcup\{\overflow\},
\]
with total operations that preserve every application defined in $D$.  The
formal representation theorem identifies separation by arbitrary intrinsic
budget-$n$ observers with separation by these finite-tag models.  Write
$x\stageeq{n}y$ when every stage-$n$ finite-tag model interprets $x$ and $y$
equally.

The original pairwise separator can be strengthened.  Instead of selecting the
subterms of two roots, select the union of the finite subterm closures of an
arbitrary finite family.  Old leaves keep their base values.  Every selected
suspended node receives the tag of its first occurrence in the selected list,
and all other behavior is sent to overflow.  A finite lookup table reproduces
each selected normal suspension exactly.  Induction over the selected closure
then shows that every selected generated Answer evaluates to its injective
encoding.

\begin{theorem}[Relative finite-pattern realization]
\label{thm:finite-pattern-realization}
Let $x_1,\ldots,x_k\in\RAlg{D}$ be any finite family.  There exist $n$ and one
compatible stage-$n$ finite-tag observer $T$ such that the canonical
interpretation
\[
 \Phi_T:\RAlg{D}\to T
\]
is injective on $\{x_1,\ldots,x_k\}$.
\end{theorem}

The construction fixes the entire base carrier pointwise; only the selected
suspended pattern consumes finite tags.

\begin{corollary}[Intrinsic finite-complement separation]
\label{thm:finite-external-separation}
For every distinct $x,y\in\RAlg{D}$ there is a compatible total extension
whose outside part has at most $|x|+|y|+1$ states and in which $x$ and $y$ have
different interpretations.  Equivalently,
\[
 \bigl(\forall n,\ x\stageeq{n}y\bigr)\Longleftrightarrow x=y.
\]
\end{corollary}

The finite-pattern theorem is structurally stronger than the pairwise bound,
while the pairwise implementation retains the sharper explicit estimate
$n=|x|+|y|$.

\begin{proposition}[Product obstruction over an infinite fixed base]
\label{prop:relative-product-obstruction}
The carrier-side condition ``pointwise-fixed base with finite outside
complement'' is not closed under ordinary binary products.  Concretely, each
embedding
\[
 \mathbb N\hookrightarrow \mathbb N\sqcup\{\star\}
\]
has a one-point outside complement, whereas the diagonal embedding of
$\mathbb N$ into
$(\mathbb N\sqcup\{\star\})^2$ has infinite outside complement.  Indeed, the
mixed points $(a,\star)$, for $a\in\mathbb N$, form an injective copy of
$\mathbb N$ outside the diagonal old image.
\end{proposition}

Thus the usual residual argument that combines finitely many pairwise
separators by taking their Cartesian product does not stay inside the
present observer class when the preserved base is infinite. The direct
finite-pattern construction in \cref{thm:finite-pattern-realization} is
therefore not replaced by that standard product trick.

\begin{remark}[What is and is not being claimed]
The selected-subterm plus sink construction has standard automata-theoretic
antecedents \cite{ComonEtAl2008}.  The contribution claimed here is not the
existence of a sink automaton in isolation.  It is the relative realization
statement for observers that must preserve a possibly infinite partial base
pointwise while budgeting only the complement.  In particular, the codomain
need not be finite when $A$ is infinite.
\end{remark}

If the base admits an injection into $\Fin(m)$, then the same observer becomes
an ordinary finite algebra with total-state bound
\[
 m+|x|+|y|+1.
\]
Hence $\RAlg{D}$ is residually finite in the ordinary sense for finitely coded
bases.  This qualitative finite-base consequence is classical in the
finite-signature ground-equational setting; the relative construction is used
because it remains meaningful when the base itself is infinite.

\begin{proposition}[Closure of the base image]
\label{prop:old-image-closed}
The image of $A$ in $\RAlg{D}$ is closed under every generated total operation
if and only if the evaluator of $D$ is total.
\end{proposition}

Thus in a genuinely partial example the preserved base is not a total
subalgebra of $\RAlg{D}$.  This is one reason ordinary finite-Rees-index
language does not literally describe the relative observer situation.

\section{The induced completion and when it is proper}
\label{sec:completion}

The finite-complement extensions of the preceding section define a descending
family of equivalence relations.  Write $x\stageeq{n}y$ when no compatible
observer using at most $n+1$ states outside the base distinguishes $x$ from
$y$.  Increasing the budget can only refine the relation, and
\cref{thm:finite-external-separation} implies
\[
  \bigcap_{n\ge 0}\stageeq{n}\;=\;=.
\]
Thus the resulting filtered space is separated.

A sequence $(x_k)$ is called Cauchy when, for every fixed $n$, all sufficiently
late terms are pairwise $\stageeq{n}$-equivalent.  Two Cauchy sequences are
identified when they eventually agree at every fixed stage.  This gives the
usual completion of a separated filtration; constant sequences embed the
original generated algebra injectively.

The question of interest is whether this embedding is surjective.  We call the
completion \emph{proper} when it contains a Cauchy class not represented by a
single generated term.

\subsection{A finite-state lemma}

Consider a sequence obtained by repeatedly applying one unary context.  At a
fixed observation stage, its interpretation is an orbit
\[
  s_{k+1}=F(s_k)
\]
of a deterministic self-map.  Suppose the part of this orbit relevant at that
stage can be encoded by $N$ states, in the sense that equality of codes at two
ordered times forces equality of the corresponding interpreted states.

\begin{lemma}[Factorial synchronization]
\label{lem:factorial-synchronization}
Under the hypothesis above, the factorial subsequence
\[
  s_{0!},s_{1!},s_{2!},\ldots
\]
is eventually constant at that observation stage.
\end{lemma}

\begin{proof}
Among the first $N+1$ coded orbit positions, two have the same code.  Hence the
orbit is eventually periodic with a period $p\le N$.  Every sufficiently large
factorial is divisible by $p$, so all sufficiently late factorial indices lie
in the same position of that eventual cycle.
\end{proof}

This is the only finite-state input used below.

\subsection{A growing orbit cannot converge to a finite term}

Assume now that
\[
  t_{k+1}=\susp(u,t_k,\old(e))
\]
for a fixed operation symbol $u$ and fixed base element $e$, and that every
step is genuinely unresolved, so the constructor size of $t_k$ strictly
increases.

Fix any generated term $x$.  Choose $N$ so large that $t_{N+1}$ is larger than
every subterm of $x$ and $t_N$.  Apply the finite-family construction of
\cref{thm:finite-pattern-realization} to the two roots $x$ and $t_N$.  The
observer remembers every selected subterm exactly.  The transition producing
$t_{N+1}$ is not among those selected transitions, so it is sent to the sink
state $\overflow$.  By construction the same unary context keeps
$\overflow$ at $\overflow$.  Therefore all later $t_k$ have the same sink
value, while $x$ retains its distinct selected value.

\begin{lemma}[No generated limit]
\label{lem:no-generated-limit}
For every generated term $x$, there is a finite-complement observer and an
index $N$ such that the observer separates $x$ from every $t_k$ with $k>N$.
Consequently no subsequence of $(t_k)$ that is eventually contained in this
tail can converge to $x$ in the filtered completion.
\end{lemma}

The argument is elementary, but it is important that the observer preserves the
entire base pointwise while using only finitely many new states.

\subsection{A sufficient criterion for properness}

Combining the two preceding lemmas yields the general mechanism used in the
examples.

\begin{theorem}[Properness from finite orbit compression]
\label{thm:compression-escape-properness}
Let $(t_k)$ be a growing orbit of the form above.  Suppose that, at every
observation stage, its interpreted trajectory admits a finite code.  Then the
factorial subsequence $(t_{k!})$ is Cauchy and has no generated limit.  Hence
the completion is proper.
\end{theorem}

\begin{proof}
At each fixed stage, \cref{lem:factorial-synchronization} makes the factorial
subsequence eventually constant, so it is Cauchy.  By
\cref{lem:no-generated-limit}, every generated candidate is separated from the
whole sufficiently late tail by some observer.  Thus no generated term
represents its Cauchy class.
\end{proof}

\subsection{Finite bases}

Suppose that $A$ injects into $\Fin(m)$ and that some application
$\eval_D(f,a,z)$ is undefined.  Starting with $c_0=\old(a)$, define
\[
  c_{k+1}=f(c_k,\old(z)).
\]
At stage $n$, the whole observer has at most $m+n+1$ states after coding the
base, so the interpreted orbit is finitely coded.  The syntax grows by one
unresolved constructor at each step.  The preceding theorem therefore gives a
new completion point.

\begin{theorem}[Finite-base criterion]
\label{prop:proper-completion}
If the base is finitely coded, the completion is proper if and only if the
original evaluator is not total.
\end{theorem}

\begin{proof}
If some base application is undefined, the construction above gives a Cauchy
sequence with no generated limit.  Conversely, if the evaluator is total, every
generated expression reduces to a base element.  Since the filtration separates
base elements, every Cauchy sequence is eventually constant and no new point is
added.
\end{proof}

\subsection{An infinite-base criterion}

The preceding argument does not require the whole base to be finite if the
chosen orbit cannot use base elements as additional memory.  Assume
\[
  \eval_D(g,a,b)=\operatorname{none}
\]
and suppose there are an operation $u$ and a base element $e$ such that
\[
  \eval_D(u,x,e)=\operatorname{some}(x)
  \qquad (x\in A).
\]
Let
\[
  c_0=\susp(g,\old(a),\old(b)),
  \qquad
  c_{k+1}=\susp(u,c_k,\old(e)).
\]
In a stage-$n$ observer, if the interpreted orbit ever enters the base, the
identity law above makes it constant from then on.  Otherwise it remains among
the $n$ tags and the overflow state.  Hence the trajectory always has a finite
code of size at most $n+2$.

\begin{theorem}[Base-fixing criterion]
\label{thm:old-fixing-proper-completion}
Under the preceding hypotheses, the completion is proper.  Moreover, for every
$n$ there are distinct iterates that cannot be separated at stage $n$; in
particular the least separation budgets along the orbit are unbounded.
\end{theorem}

\begin{proof}
The properness statement follows from
\cref{thm:compression-escape-properness}.  For the quantitative statement, the
finite code has at most $n+1$ non-base states unless the orbit has already
stabilized in the base.  Factorial synchronization therefore gives
\[
  c_{(n+1)!}\stageeq{n}c_{(n+2)!}.
\]
The two terms are syntactically distinct because each iteration adds one
constructor.
\end{proof}

The finite-base and base-fixing hypotheses are independent.  A finite partial
algebra need not have a unary context fixing every base element, while natural
arithmetic has the context $x\mapsto x+0$ but has an infinite base.

\subsection{Comparison with the prefix-depth metric}

The filtration above is not merely the usual prefix metric on trees.  Let
\[
  b_0=\old(a),\qquad b_{k+1}=\susp(f,b_k,\old(z)),
\]
with $\eval_D(f,a,z)$ undefined.  For every fixed syntactic depth, the sequence
$(b_k)$ is eventually constant on that prefix and is therefore Cauchy in the
standard tree metric.

By contrast, a single one-tag finite-complement observer can be chosen so that
the comb context alternates between its tag and the overflow state.  The same
observer then distinguishes every adjacent pair $b_k,b_{k+1}$, regardless of
how deep the new constructor occurs.

\begin{proposition}[The two Cauchy notions differ]
\label{prop:depth-observation-contrast}
There exists a generated comb sequence that is Cauchy for prefix-depth
agreement but is not Cauchy for the finite-complement filtration.
\end{proposition}

Thus the completion studied here is not simply the ordinary infinite-tree
completion written in different notation.

\section{Exact Equational Conservativity}\label{sec:equations}

The completion adds observational limit points, but it does not create new
universal equations in the language containing variables, named base constants,
and the primitive binary operations.

Term evaluation commutes with the canonical embedding
$\RAlg{D}\to\Comp{\RAlg{D}}$.  Moreover, every term operation preserves each
stage relation because the primitive operations do.  Given a completed
assignment and a finite stage $n$, stage-density supplies generated
approximants for all variables.  If an equation holds for all generated
assignments, the two completed term values are therefore equivalent at every
finite stage.  Separatedness upgrades agreement at all stages to equality.
The converse follows immediately by restricting a completed equation to
embedded generated assignments.

\begin{theorem}[Exact equation preservation]
\label{thm:equational-conservativity}
For every constant-bearing algebraic equation $s=t$,
\[
  \RAlg{D}\models s=t
  \quad\Longleftrightarrow\quad
  \Comp{\RAlg{D}}\models s=t.
\]
\end{theorem}

\begin{remark}[Scope]
This is a theorem about universal equations.  Arbitrary quasi-equations do not
follow from density alone; conditional laws require additional hypotheses on
the premises.  No unrestricted quasi-equational conservativity claim is made.
\end{remark}

\section{Arithmetic Examples}\label{sec:arithmetic}

The construction is not tied to arithmetic, but natural and integer examples
make the arbitrary-base properness mechanism concrete.  In both cases division
is partial at denominator zero, while addition supplies the base-fixing context
$x\mapsto x+0$.

\subsection{Natural numbers}

Take $+$, $\cdot$, and $/$ on $\mathbb N$, with division undefined exactly at
zero denominator.  For $b\ne0$, resolution returns the embedded natural
$a/b$.  At zero,
\[
  a/0=\susp(\mathsf{div},\old(a),\old(0)).
\]
Distinct numerators give distinct singular Answers, and no singular Answer is
an embedded natural.  Lean exposes these statements as
\leanname{ResolutionSemantics.NatDivision.singularFamilyInjective} and
\leanname{ResolutionSemantics.NatDivision.singularFamilyDisjoint}.

The base carrier is infinite and admits no finite injection into $\Fin(m)$.
Nevertheless, addition satisfies $a+0=a$ for every base natural.  Starting from
the undefined seed $0/0$ and iterating the context $x\mapsto x+0$ gives the
old-fixing orbit of \cref{thm:old-fixing-proper-completion}.  Hence its
factorial sample is Cauchy with no generated limit, the observational
completion is proper, and finite-tag separation ranks are unbounded.  The
criterion is exposed by
\leanname{ResolutionSemantics.NatDivision.oldFixingCriterion} and the rank
consequence by \leanname{ResolutionSemantics.NatDivision.separationRanksUnbounded}.

This example also illustrates why the filtration is not merely syntactic depth.
A different iterated-division comb can be separated already at stage zero by
using unused base naturals as semantic memory, while the repeated-$+0$ orbit
cannot: preservation forces every base natural to be fixed by the context.
Thus the relative observer class can see deep syntax aggressively in one family
and compress it strongly in another.

\subsection{Integers}

The integer instance additionally includes subtraction.  Zero-denominator
applications again remain structured suspended Answers rather than ordinary
integers.  The singular family is injective and completed $0/0$ remains outside
the embedded integer image; see
\leanname{ResolutionSemantics.IntDivision.singularFamilyInjective} and
\leanname{ResolutionSemantics.IntDivision.zeroDivZeroNotOld}.

Because $a+0=a$ also holds over $\mathbb Z$, the same compression--escape
mechanism proves properness and unbounded separation rank.  The relevant public
declarations include
\leanname{ResolutionSemantics.IntDivision.completionEmbeddingNotSurjective} and
\leanname{ResolutionSemantics.IntDivision.separationRanksUnbounded}.

\begin{resultbox}
Resolution does not assign an ordinary numerical value to $0/0$.  It preserves
its unresolved provenance as a structured Answer in a conservative total
extension.
\end{resultbox}

This differs from meadows, common meadows, and wheels, where singular
expressions are totalized by algebraic conventions or designated exceptional
values \cite{BergstraHirshfeldTucker2009,BergstraPonseCommonMeadows,Carlstrom2004}.
The goal here is not to replace those systems, but to retain the computation
that failed and let compatible observers decide how much of that structure to
remember.

\section{A Lean-Checked Bridge to Diaconescu's Encoding}
\label{sec:diaconescu-bridge}

The main separation and completion theorems above are one-sorted and binary.
Their free-completion substrate sits inside a broader partial-to-total theory.
To make that provenance explicit, the repository contains an independent
many-sorted Lean development of the operations-only fragment of Diaconescu's
encoding \cite{Diaconescu2009}.

A many-sorted signature has sorts $S$, designated-total operation symbols
$TF$, and partial operation symbols $PF$, each with arbitrary finite arity and
typed argument and result sorts.  Encoded total models extend the data carriers
with an auxiliary truth carrier and existence-equality operations.  The Horn
conditions $\Gamma$ ensure that the defined elements form the partial reduct
$\beta(M)$.

The translation $\alpha$ sends existence-equality atoms to equality with the
truth constant and guards universal quantification by definedness.  Lean proves
the term soundness and completeness lemmas and hence the satisfaction
condition
\[
  M\models\alpha(\rho)
  \quad\Longleftrightarrow\quad
  \beta(M)\models\rho.
\]
The relevant declarations are
\leanname{Resolution.Diaconescu.ManySorted.term\_evalPartial\_sound},
\leanname{Resolution.Diaconescu.ManySorted.term\_evalPartial\_complete}, and
\leanname{Resolution.Diaconescu.ManySorted.alpha\_satisfaction\_condition}.

For every partial algebra $D$, the canonical model $\gamma(D)$ uses normalized
formal applications as its total carrier.  Lean checks $\gamma(D)\models\Gamma$
and a canonical partial-algebra equivalence
$\beta(\gamma(D))\simeq D$ through
\leanname{Resolution.Diaconescu.ManySorted.gamma\_satisfiesGamma} and
\leanname{Resolution.Diaconescu.ManySorted.betaGammaPartialAlgEquiv}.
The universal lift from a partial homomorphism into a reduct is formalized as a
hom-set equivalence; identity, composition, naturality, unit/counit behavior,
and both triangle identities are checked.

At fixed signature and universe bound, the development proves the initiality
transfers needed for the operations-only quasi-existence theories considered
here.  A conditional term-model construction supplies the initial translated
encoded model, and the $\beta$ reduct yields the corresponding initial partial
model.  Semantic consequence is preserved and reflected by the translation.
Representative declarations include
\leanname{Resolution.Diaconescu.ManySorted.gammaBetaHomEquiv},
\leanname{Resolution.Diaconescu.ManySorted.gamma\_preserves\_initiality},
\leanname{Resolution.Diaconescu.ManySorted.beta\_of\_initial\_encoded\_is\_initial},
\leanname{Resolution.Diaconescu.ManySorted.semantic\_consequence\_equivalence},
and
\leanname{Resolution.Diaconescu.InitialExistence.quasiTheory\_has\_initial\_partial\_model}.

Finally, when the signature has one sort, no designated-total symbols, and only
binary partial symbols, Lean constructs canonical equivalences between the
many-sorted generated carrier and the original binary Resolution carrier, and
between the corresponding encoded models.  The principal declarations are
\leanname{Resolution.Diaconescu.BinarySpecialization.generatedEquiv} and
\leanname{Resolution.Diaconescu.BinarySpecialization.gammaEquiv}.

\begin{remark}[Role of the bridge]
The bridge verifies that the binary implementation used by the observer theory
is a specialization of an established many-sorted partial-to-total encoding.
It is supporting formal provenance.  We do not claim Diaconescu's encoding,
the adjunction, or the existence of free partial completions as new results.
Nor does the present paper extend the finite-complement separation and
compression--escape completion theory to the full many-sorted setting; that
remains an open direction.
\end{remark}

\section{Relation to Prior Work}\label{sec:related-work}

\subsection{Free completion and partial algebras}
The free compatible completion underlying \cref{sec:free-completion} belongs to
an established line of work on partial algebras and Schmidt-type completions
\cite{ClimentVidalCosmeLlopez2023}.  Diaconescu's many-sorted encoding provides
a closely related partial-to-total perspective \cite{Diaconescu2009}.  These
results motivate the substrate but are not novelty claims of the present paper.

\subsection{Recognizable trees and finite separators}
For finite bases and finite signatures, the normalized ground-equational
presentation lies within classical tree-automata and recognizability theory
\cite{Kozen1977,Brainerd1969,DauchetTison1990,ComonEtAl2008}.  Finite selected
subterm closures plus a sink state are familiar constructions.  The relative
feature studied here is the requirement that every observer preserve an
arbitrary partial base pointwise while only the complement over that base is
budgeted.  This permits infinite codomains with finite external complexity.

There is a concrete obstruction to reducing the simultaneous finite-pattern
statement to the usual ``separate pairs and take their product'' argument.
The obstruction is already visible at carrier level.  For an infinite base
$A$, each one-point extension $A\hookrightarrow A\sqcup\{\star\}$ has finite
relative complement.  But the Cartesian product of two such extensions, with
$A$ embedded diagonally, contains all mixed points $(a,\star)$ outside the old
image.  Hence its complement is infinite.  This failure of product closure is
Lean-checked for $A=\mathbb N$ in
\texttt{ResolutionRelativeProductObstruction.lean}.  Thus abstract
finite-family discrimination may be classical, but the standard product
construction does not establish it inside the restricted observer class used
here when the preserved base is infinite; the direct finite-pattern
construction is doing additional work.

\subsection{Filtered and profinite completions}
Cauchy completion of separated filtrations and factorial synchronization of
finite deterministic dynamics are standard.  Profinite completions organize
observations by finite quotients or finite recognizing algebras
\cite{DaveyGouveiaHaviarPriestley2011,ChenAdamekMiliusUrbat2016,Moreau2024}.
Resolution stages differ because the preserved base may be infinite.  The
compression theorem deliberately isolates the standard dynamic ingredient:
only the trajectory relevant to the chosen unary family needs a finite code.
The new structural interaction is with the relative finite-pattern escape
observer that prevents convergence to any finite generated candidate.

\subsection{Finite complements and Rees-index analogies}
Finite-complement extensions and residual finiteness are classical in semigroup
theory \cite{CainMaltcev2013,RuskucThomas1998}.  The analogy is useful but not
literal here.  In a genuinely partial Resolution algebra, the embedded old
carrier is not closed under the generated total operations; by
\cref{prop:old-image-closed}, closure occurs exactly when the evaluator was
already total.  Hence the observer class is better described as relative to a
pointwise-fixed partial base than as an ordinary finite-Rees-index extension.

\subsection{Tree metrics}
Classical infinite-tree completions are based on prefix depth
\cite{ArnoldNivat1980,Courcelle1983}.  Resolution uses semantic finite-tag
observation instead.  The formal comb counterexample in \cref{sec:completion}
shows that the two Cauchy notions differ: a sequence may be depth-Cauchy while a
fixed low-budget Resolution observer continues to distinguish adjacent terms.

\subsection{Contribution and scope}
The contribution claimed by this paper is concentrated in the following
package:
\begin{enumerate}
  \item an intrinsic finite-complement observer class over a pointwise-fixed,
        possibly infinite partial base, together with finite-pattern
        realization;
  \item the induced observational filtration and its trajectory-level
        compression theorem;
  \item the Resolution-specific finite-pattern escape mechanism and the
        resulting compression--escape properness theorem;
  \item the exact finite-base criterion, infinite-base old-fixing instance,
        sharper unbounded-rank witnesses, and exact equational conservativity.
\end{enumerate}
The factorial-periodicity argument itself is treated as standard machinery.
Likewise, the many-sorted Diaconescu development is a verified compatibility and
provenance layer rather than an independent novelty claim.

\section{Machine-Checked Companion}\label{app:formalization}

\noindent
Nothing in this appendix is used by the mathematics above.  Every definition,
theorem and proof in the body is complete in ordinary mathematical language and
should be read and judged as such.  The development recorded here is a
correctness check on the author's own reasoning, not evidence offered to the
reader, and it establishes neither significance nor priority.  It is reported
because it exists, and because its trust boundary should be stated explicitly
rather than left implicit.

The repository pins Lean~4.33.0.  The proof package contains a small
publication-facing facade in \texttt{ResolutionMasterTheorems.lean} that
exposes the structural results used in the reorganized paper.  In particular,
the following declarations are checked directly:
\begin{lstlisting}[style=lean]
#check ResolutionSemantics.MasterTheorems.relativeFinitePatternRealization
#check ResolutionSemantics.MasterTheorems.orbitCompressionCauchy
#check ResolutionSemantics.MasterTheorems.finitePatternEscapeNoGeneratedLimit
#check ResolutionSemantics.MasterTheorems.compressedEscapingOrbitNotComplete
#check ResolutionSemantics.MasterTheorems.compressedEscapingOrbitEmbeddingNotSurjective
#check ResolutionSemantics.MasterTheorems.finiteBasePropernessCriterion
#check ResolutionSemantics.MasterTheorems.oldFixingPropernessViaCombinedMaster
#check ResolutionSemantics.MasterTheorems.oldFixingRanksUnbounded
\end{lstlisting}

The finite-pattern result is implemented in
\texttt{ResolutionFinitePatternRealization.lean}.  The trajectory-compression
abstraction is in \texttt{ResolutionOrbitCompressionMaster.lean}; its finite
code applies only to the relevant deterministic orbit and does not require a
finite ambient observer carrier.  The anti-limit theorem is implemented in
\texttt{ResolutionFinitePatternAntiLimit.lean}; it formalizes the selected
candidate/prefix construction, overflow entry, and overflow absorption.  Both
the finite-base and old-fixing properness theorems are re-derived as instances
of the combined theorem.

The command
\begin{lstlisting}[style=lean]
bash scripts/verify.sh
\end{lstlisting}
checks the manuscript-linked public API, compiles the proof dependency chain,
rejects direct proof holes or local axioms in the paper sources, and runs the
axiom audit.  The accepted theorem dependencies are restricted to
\texttt{propext}, \texttt{Classical.choice}, and \texttt{Quot.sound}.

The formalization uses classical finite indexing and witness selection in
explicit places.  The resulting trust claim is therefore relative to Lean's
kernel, the imported standard library, and the ordinary classical principles
reported by the audit.  No foundation-independent consistency claim is made.

\subsection{Theorem map}

The declarations corresponding to the numbered results of the body are
listed below, grouped by section.  They are provided for traceability only.

\paragraph{Free completion substrate (\S
ef{sec:free-completion})}
\begin{itemize}\setlength{\itemsep}{0pt}
  \item \leanname{ResolutionSemantics.freeCompletionUniversal}
  \item \leanname{ResolutionSemantics.expressionKernelQuotientInjective}
  \item \leanname{ResolutionSemantics.expressionKernelQuotientSurjective}
\end{itemize}

\paragraph{Observers and finite-pattern realization (\S
ef{sec:finite-separation})}
\begin{itemize}\setlength{\itemsep}{0pt}
  \item \leanname{ResolutionSemantics.MasterTheorems.relativeFinitePatternRealization}
  \item \leanname{ResolutionSemantics.finiteSeparationBound}
  \item \leanname{ResolutionSemantics.MasterTheorems.relativeFiniteComplementSeparation}
  \item \leanname{ResolutionSemantics.RelativeObservers.finiteRelativeComplement_not_closed_under_binary_products}
\end{itemize}

\paragraph{Completion, compression, escape (\S
ef{sec:completion})}
\begin{itemize}\setlength{\itemsep}{0pt}
  \item \leanname{ResolutionSemantics.MasterTheorems.orbitCompressionCauchy}
  \item \leanname{ResolutionSemantics.MasterTheorems.finitePatternEscapeNoGeneratedLimit}
  \item \leanname{ResolutionSemantics.MasterTheorems.compressedEscapingOrbitEmbeddingNotSurjective}
  \item \leanname{ResolutionSemantics.MasterTheorems.compressedEscapingOrbitNotComplete}
  \item \leanname{ResolutionSemantics.MasterTheorems.finiteBasePropernessCriterion}
  \item \leanname{ResolutionSemantics.MasterTheorems.oldFixingPropernessViaCombinedMaster}
  \item \leanname{ResolutionSemantics.MasterTheorems.oldFixingRanksUnbounded}
  \item \leanname{ResolutionSemantics.completionComplete}
  \item \leanname{Resolution.External.FilteredTotalAlg.completedResolutionFilteredAlgebra\_initial}
  \item \leanname{Resolution.Probe.depthCauchy\_not\_observationalCauchy}
  \item \leanname{Resolution.Probe.comb\_separatesAt\_one}
  \item \leanname{Resolution.Probe.NatProbe.natComb\_pairwise\_separated\_at\_zero}
\end{itemize}

\paragraph{Equational conservativity (\S
ef{sec:equations})}
\begin{itemize}\setlength{\itemsep}{0pt}
  \item \leanname{ResolutionSemantics.equationConservative}
\end{itemize}

\paragraph{Arithmetic instances (\S
ef{sec:arithmetic})}
\begin{itemize}\setlength{\itemsep}{0pt}
  \item \leanname{ResolutionSemantics.NatDivision.singularFamilyInjective}
  \item \leanname{ResolutionSemantics.NatDivision.singularFamilyDisjoint}
  \item \leanname{ResolutionSemantics.NatDivision.oldFixingCriterion}
  \item \leanname{ResolutionSemantics.NatDivision.separationRanksUnbounded}
  \item \leanname{ResolutionSemantics.IntDivision.singularFamilyInjective}
  \item \leanname{ResolutionSemantics.IntDivision.zeroDivZeroNotOld}
  \item \leanname{ResolutionSemantics.IntDivision.completionEmbeddingNotSurjective}
  \item \leanname{ResolutionSemantics.IntDivision.separationRanksUnbounded}
\end{itemize}

\paragraph{Many-sorted localization (\S
ef{sec:diaconescu-bridge})}
\begin{itemize}\setlength{\itemsep}{0pt}
  \item \leanname{Resolution.Diaconescu.ManySorted.term\_evalPartial\_sound}
  \item \leanname{Resolution.Diaconescu.ManySorted.term\_evalPartial\_complete}
  \item \leanname{Resolution.Diaconescu.ManySorted.alpha\_satisfaction\_condition}
  \item \leanname{Resolution.Diaconescu.ManySorted.gamma\_satisfiesGamma}
  \item \leanname{Resolution.Diaconescu.ManySorted.betaGammaPartialAlgEquiv}
  \item \leanname{Resolution.Diaconescu.ManySorted.gammaBetaHomEquiv}
  \item \leanname{Resolution.Diaconescu.ManySorted.gamma\_preserves\_initiality}
  \item \leanname{Resolution.Diaconescu.ManySorted.beta\_of\_initial\_encoded\_is\_initial}
  \item \leanname{Resolution.Diaconescu.ManySorted.semantic\_consequence\_equivalence}
  \item \leanname{Resolution.Diaconescu.InitialExistence.quasiTheory\_has\_initial\_partial\_model}
  \item \leanname{Resolution.Diaconescu.BinarySpecialization.generatedEquiv}
  \item \leanname{Resolution.Diaconescu.BinarySpecialization.gammaEquiv}
\end{itemize}

\section{Limitations and Open Questions}\label{sec:limitations}

The structural reduction clarifies several boundaries rather than eliminating
them.

\begin{enumerate}
  \item The finite-pattern realization theorem is proved for the main
        one-sorted binary Resolution kernel.  An independent Lean module checks
        the same realization mechanism for arbitrary finite arities on finite
        child-closed normalized raw patterns, but the full observational
        completion theory has not yet been lifted to the many-sorted setting.
  \item The pairwise constructor-size tag bound is a general upper bound, not
        an optimal state-complexity theorem.  Determining the optimal relative
        separation rank as a function of Answer size and signature/base
        hypotheses remains open.
  \item The finite-base theorem gives an exact properness classification for
        finitely coded carriers.  The old-fixing condition is a sufficient
        infinite-base criterion, not a necessary one.  A full arbitrary-carrier
        characterization of properness remains open.
  \item Trajectory compression is deliberately weak: it constrains only the
        observer orbit relevant to a selected generated family.  Other
        trajectories in the same observer may have unbounded or genuinely
        infinite complexity.  This is a feature of the theorem, but it limits
        conclusions about global finite-state structure.
  \item Finite-pattern escape is currently formalized for a fixed right unary
        syntactic context $t\mapsto\susp(u,t,\old(e))$.  More general one-hole
        contexts and arbitrary finitary contexts are natural extensions.
  \item The many-sorted bridge checks the concrete operations-only fragment,
        fixed-signature adjunction laws, semantic consequence, and the
        initial-model construction used here.  It does not formalize the full
        surrounding institution/comorphism infrastructure or arbitrary
        signature morphisms from Diaconescu's broader theory.
  \item Exact conservativity is proved for universal equations with named base
        constants.  Arbitrary conditional equations require additional
        hypotheses.
  \item The construction is semantic, not an efficiency claim.  Exact
        provenance retention can produce large finite Answers and large
        candidate-tailored observer tables.
\end{enumerate}

\begin{openquestion}[Relative compression--escape criteria]
Can the combined properness theorem be characterized directly in terms of the
partial algebra $D$, without first exhibiting a particular compressed escaping
unary orbit?  Such a criterion would extend the exact finite-base theorem
beyond finitely coded carriers.
\end{openquestion}

\begin{openquestion}[General contexts]
How far does finite-pattern escape extend from the fixed-right unary context to
arbitrary finite one-hole contexts, many-sorted signatures, or finitely
branching context systems?  The existing finite-pattern realization theorem
suggests that the anti-limit mechanism is more general than the current unary
formalization.
\end{openquestion}

\section{Conclusion}

Resolution semantics starts from a familiar free-compatible-completion idea:
defined base applications reduce, while undefined applications remain
structured finite Answers.  The present paper's contribution lies in what is
built over that substrate.

The first structural principle is relative finite-pattern realization.  A
finite family of generated Answers can be realized simultaneously inside one
compatible observer that fixes the entire partial base pointwise and uses only
finitely many new states beyond it.  Pairwise finite-complement separation is
the two-point consequence, and finite-tag models provide canonical stage
representatives.  These stages induce a separated observational filtration.

The second part of the story explains why its completion can be genuinely
larger than the generated Answer algebra.  Finite trajectory compression makes
factorial samples Cauchy even when the ambient observer carrier is infinite.
Finite-pattern escape supplies the opposite force: every finite generated
candidate can be retained exactly while a sufficiently late continuation of a
growing orbit is forced into absorbing overflow.  Thus the same relative
observer technology both compresses infinite behavior at every fixed stage and
exposes escape beyond every finite candidate.  Their interaction forces new
completion points.

This compression--escape theorem unifies the finite-base and infinite-base
old-fixing properness results.  For finitely coded bases, properness is exactly
equivalent to partiality of the evaluator.  For arbitrary bases, a pointwise
base-fixing unary context supplies a compressed trajectory; the sharpened
stage-$n$ witness $c_{(n+1)!}\stageeq{n}c_{(n+2)!}$ yields unbounded separation
rank.  Natural and integer arithmetic instantiate this mechanism through the
context $x\mapsto x+0$, while singular division remains structured rather than
being assigned an ordinary numerical value.

The observational completion carries canonical extensions of the primitive
operations and preserves exactly the universal equations of the generated
algebra in the language with named base constants.  Its filtration is not the
classical prefix-depth filtration on trees: a fixed one-tag observer can
separate adjacent unresolved combs at arbitrary depth, and in natural
arithmetic the retained base supports an even sharper zero-tag separator.
Finite semantic observers can therefore compress or distinguish arbitrarily
deep syntax according to compatibility with the preserved partial base.

The many-sorted Lean development provides a separate provenance result.  It
checks the relevant $\Gamma$--$\alpha$--$\beta$--$\gamma$ encoding machinery
and proves that the original binary Resolution core is the expected
single-sorted specialization.  This distinguishes inherited partial-to-total
infrastructure from the relative observation and completion results developed
here.

The remaining questions now have a cleaner form: determine optimal relative
separation complexity, characterize properness over arbitrary infinite bases,
and extend finite-pattern escape and trajectory compression to general finitary
and many-sorted contexts.  The current formalization establishes the binary
compression--escape architecture as a verified starting point for those
extensions.


\begingroup
\small
\setlength{\bibsep}{0.25em}
\bibliographystyle{abbrvnat}
\bibliography{../references,../references-completion,../references-separation}
\endgroup

\end{document}
